\chapter{Conclusion}

The Thesis and Project Management System (TPMS) was designed, developed, and validated as a comprehensive, role-aware, production-grade web platform for digitalizing and automating the complete administrative and evaluative lifecycle of Bachelor Minor/Major projects and Master research theses at the Department of Electronics and Computer Engineering, Pulchowk Campus, Institute of Engineering, Tribhuvan University.

The system successfully addresses all six primary problem areas identified in the problem formulation: decentralized record-keeping, lack of multi-project engagement controls, supervisory conflict-of-interest vulnerabilities, manual mark tabulation, cumbersome official documentation, and absence of immutable audit logging.

Key technical contributions of the project include:
\begin{itemize}
    \item A strict \textbf{program-scoped degree isolation} architecture supporting six academic programs across Bachelor and Master tiers from a single unified deployment.
    \item An \textbf{automated Engagement Guard} with database-level and application-layer enforcement preventing dual student enrollment across project groups and theses.
    \item An \textbf{algorithmic Conflict-of-Interest Filter} that dynamically excludes project supervisors from examiner selection dropdowns, eliminating ethical evaluation conflicts.
    \item A \textbf{multi-tier evaluation rubric engine} implementing the DOECE grading schemes for Minor (50 marks), Major (100 marks), and Master Thesis (300 marks) evaluations.
    \item A \textbf{server-side headless browser PDF pipeline} using Puppeteer for instant generation of official TU Pulchowk Campus A4 evaluation sheets with number-to-words mark conversion, eliminating hours of manual typesetting.
    \item A \textbf{coordinator response matrix} with group deduplication and 1-click fuzzy supervisor matching, reducing per-group supervisor allocation time from 15 minutes to under 30 seconds.
    \item A \textbf{program-scoped immutable audit trail} providing coordinators and maintainers with complete operational traceability.
\end{itemize}

Functional validation through unit testing (36 test cases, 100\% pass rate), integration testing (4 end-to-end workflows), security testing (5 attack vector scenarios), and User Acceptance Testing (4.47/5.0 average satisfaction) confirmed that the system meets and exceeds all defined functional, performance, and usability requirements.

The successful implementation of TPMS demonstrates that a full-stack Node.js/React/PostgreSQL architecture with carefully engineered domain-specific business logic can effectively replace manual, paper-based academic administration workflows in institutional engineering education. The platform is ready for intranet deployment at DOECE, Pulchowk Campus, and provides a replicable, extensible model for academic project management digitalization that can be adapted by other departments and engineering institutions with analogous administrative structures.
